% Lynceus: Cost-Model-Driven Query Routing for Heterogeneous GPU-CPU Systems
% Final reconstructed paper (placeholders fully resolved), main body Sections 1-7.
% Counterpart of des_loc_reconstructed.tex; appendix excluded.
% Claude-6 (M126-M150): full readable reconstruction from the templated
%   lynceus_reconstructed_EN.tex (Claude-1..5) into real academic prose.
%   Sections: 1 Introduction, 2 Lynceus Algorithm, 3 Cost-Model Correctness,
%             4 Experimental Design, 5 Evaluation (RQ1-RQ6), 6 Related Work,
%             7 Conclusion.

\documentclass{article}
% Fall back to article geometry if the NeurIPS style is unavailable.
\IfFileExists{neurips_2026.sty}{\usepackage{neurips_2026}}{\usepackage[margin=1in]{geometry}}
\usepackage{amsmath,amssymb,amsfonts}
\usepackage{amsthm}
\usepackage{algorithmic}
\usepackage{algorithm}
\usepackage{graphicx}
\usepackage{booktabs}
\usepackage{multirow}
\usepackage{hyperref}

\newtheorem{theorem}{Theorem}
\newtheorem{proposition}{Proposition}
\newtheorem{assumption}{Assumption}

\title{Lynceus: Cost-Model-Driven Query Routing for Heterogeneous GPU-CPU Systems}

\author{
  Dylan Yunlon$^{\dagger,1}$ \\
  $^{1}$Independent Research \\
  $^{\dagger}$\texttt{dogechat@163.com}
}

\begin{document}
\maketitle

\begin{abstract}
Analytical query processing increasingly runs on heterogeneous nodes that pair
many-core CPUs with GPUs, yet most systems route every query to a single fixed
target. Routing whole queries to one device ignores that each relational operator
has a very different cost profile on the CPU and the GPU, and that moving data
across the PCIe/NVLink fabric into device memory incurs a transfer cost that is
frequently omitted from estimation. Optimistic, transfer-blind cost models thus
mispredict the cheaper device and overload one side of the node. We present
Lynceus, a system that drives per-operator device assignment from a single,
transfer-aware cost model: it estimates an itemized cost on every device, prices
cross-device data movement as a multi-hop shortest path over the hardware
topology, and routes each operator to its lowest-cost target. Lynceus prices
virtual indexes so that index-eligible queries pay a traversal rather than a full
scan, keeps hot data resident through a block-level index cache, and compresses
its statistics with block-scaled FP8 so that cost tables occupy a fraction of the
memory footprint. We give a set of binding correctness contracts---a single
query's predicted latency equals its full critical path including transfer; per
query estimation never invents cross-device overlap; logical-table identity is an
explicit field; and fabric cost follows multi-hop shortest paths---and we show
that the only legitimate source of pipeline speedup is batching independent
queries, governed by the Megatron-LM pipeline-bubble fraction. On a TPC-H SF100
workload over a multi-socket CPU paired with multiple GPUs, Lynceus achieves
$1.3$--$2.1\times$ end-to-end latency reductions over single-device and static,
transfer-blind routing, remains robust when data resides on a second NUMA node
reachable only multi-hop, and preserves routing quality when its cost tables are
quantized to block-scaled E4M3.
\end{abstract}

%=============================================================================
\section{Introduction}
\label{sec:introduction}
%=============================================================================

Modern analytical data systems increasingly run on heterogeneous nodes, pairing
many-core CPUs with one or more GPUs to improve throughput and to fit larger
working sets in aggregate device memory. However, routing every query to a fixed
execution target---whether GPU-only or CPU-only---ignores two facts that dominate
performance on such hardware. First, each relational operator (scan, filter,
join, aggregate, sort) has a sharply different cost profile on the CPU than on
the GPU. Second, running an operator on the GPU requires moving its input across
the PCIe/NVLink fabric into device memory, and this transfer cost is substantial
and frequently omitted. Early systems side-stepped both issues with hand-written
static hybrid routing, dispatching whole queries to one device according to a
fixed rule rather than estimating cost per device. Modern heterogeneous
workloads, such as mixed OLAP query streams, exhibit operator-level cost
asymmetry that such coarse routing cannot exploit.

Some approaches extend rule-based query optimization to the heterogeneous setting
by costing whole-query plans on a single device, which poses three challenges.
First, they omit the cross-device transfer cost, treating data as already
resident and producing optimistic estimates that favor the GPU even when the
load into device memory makes it the slower choice. Second, charging only
on-device work accumulates systematic error and provides no principled mechanism
to keep hot data resident in faster memory, leaving such methods brittle in
bandwidth-constrained deployments. Third, ignoring multi-hop paths through the
device fabric destabilizes routing on multi-socket machines, where data on a
second NUMA node reaches a GPU only through an intermediate hop.

Static hybrid routing addresses part of this picture: pinning a query to the
device its bulk operator favors can beat naive single-device routing and remains
robust to adding new devices. However, it requires transfer cost to be folded
into compute and never priced separately, which hides cross-device data-movement
overhead relative to a transfer-aware estimate and prevents the planner from
choosing the cheaper device per operator. Hence, we aim to answer the following
question:

\begin{center}
\emph{Can independently pricing operator execution and cross-device data transfer
improve data-movement efficiency for heterogeneous query processing while
preserving routing correctness and robustness?}
\end{center}

As a result of our inquiry, we propose Lynceus, a system that drives per-operator
device assignment from a single, unified, heterogeneity-aware cost model. By
estimating an itemized cost on every device and routing each operator to its
lowest-cost target, Lynceus reduces cross-device data movement by keeping hot
data resident and moving it only when it pays off. The cost model decomposes a
query's cost into input/output, compute, transfer, index, and sort terms on a
common microsecond scale, so the GPU's parallelism is weighed directly against
the fixed cost of loading data on-device.

\paragraph{Contributions.}
\begin{enumerate}
    \item \textbf{Correct cost modeling.} We give a set of binding cost-model
    contracts for Lynceus (Section~\ref{sec:guarantees}): a single query's
    predicted latency is its full critical path including transfer; per-query
    estimation never invents cross-device overlap; logical-table identity is an
    explicit field rather than string-inferred; and fabric transfer follows
    multi-hop shortest paths. Correct transfer accounting is what makes
    per-device estimates trustworthy.

    \item \textbf{Data-movement reduction.} We show empirically that
    data-resident routing and a lightweight block-level index cache reduce
    cross-device data movement by skipping redundant transfers, yielding
    end-to-end latency reductions over single-device routing on heterogeneous
    hardware.

    \item \textbf{Scalability across workloads.} We validate Lynceus on a
    standard analytical benchmark (TPC-H at scale factor 100), demonstrating
    competitive latency against both static hybrid and single-device routing
    while routing each operator to the cheaper device.

    \item \textbf{Hardware robustness.} Unlike fixed-rule routing, Lynceus
    reasons over the full hardware topology and prices block-scaled FP8
    statistics, enabling it to exploit second-NUMA-node memory and to operate on
    compact cost tables under tight device-memory budgets.
\end{enumerate}

%=============================================================================
\section{Lynceus: Cost-Model-Driven Routing}
\label{sec:lynceus}
%=============================================================================

\subsection{Per-Device Cost and the Role of Data Placement}
\label{sec:cost}

We begin by characterizing the relationship between where data resides and the
cost of running an operator on each device, and how this can be leveraged to
lower cross-device data movement. Lynceus models the cost of routing a query
$q$ to device $s$ as a sum of itemized terms,
\begin{equation}
\label{eq:cost}
\mathrm{cost}(q, s) = c_{\mathrm{io}} + c_{\mathrm{compute}} + c_{\mathrm{transfer}} + c_{\mathrm{index}} + c_{\mathrm{sort}},
\end{equation}
all expressed in microseconds so that devices with very different execution
models can be compared on a common scale. For single-device routing the choice
is correct only if $c_{\mathrm{transfer}}$ is counted: a GPU may win on
$c_{\mathrm{compute}}$ yet lose once the load of its input across the fabric into
device memory is added. Conversely, when the input is already resident the
transfer term vanishes and the GPU path dominates on far more operators.

A useful summary is the per-device cost shape. A CPU prices a sequential scan by
$c_{\mathrm{io}} \propto \text{rows} \times \varepsilon_{\mathrm{scan}}$, where
$\varepsilon_{\mathrm{scan}}$ is a per-row coefficient, while a GPU amortizes the
same scan over massive parallelism but pays a fixed transfer to bring the data
on-device. Using these coefficients, a sequential CPU scan and a
transfer-plus-parallel GPU scan cross over at a workload-dependent cardinality,
which the cost model locates per query. Below the crossover the CPU is cheaper;
above it the GPU wins despite the transfer.

Because the input may sit on a remote NUMA node, cross-device transfer is not a
single edge but the best path through the device fabric. On a dual-socket box the
second NUMA node may have no direct GPU link, so data reaches a GPU via a
multi-hop route (for example, remote socket $\to$ local socket $\to$ GPU). The
fabric cost is therefore the minimum-cost path, additive over hops, computed as a
single-source shortest path over the topology graph with non-negative edge
weights; a genuinely disconnected pair is priced as unreachable. From this it
follows that the resident location of data should inform routing: if hot data
already occupies device memory the transfer term drops out and routing the
operator there is cheap, whereas cold data may be cheaper to process in place on
the CPU.

Where the static estimate is biased---for example, because cardinality was
misestimated---an online policy corrects it. Lynceus tracks a per-device
bias-correction ratio between observed and estimated latency, updates it with an
exponential moving average, $\mathrm{bias} \gets \alpha \cdot \mathrm{ratio} +
(1-\alpha)\,\mathrm{bias}$, and shifts subsequent dispatch decisions through a
load-balancing margin. The decay rate $\alpha$ sets the memory half-life of the
correction, damping both transient queueing stalls and needless high-frequency
re-probing that, in the extreme, could exhaust device memory.

\subsection{The Lynceus Router}
\label{sec:router}

Algorithm~\ref{alg:lynceus} summarizes routing. For each query Lynceus estimates
the itemized cost of Equation~\eqref{eq:cost} on every device, optionally pricing
a virtual index so that an index-eligible query pays a B-tree traversal rather
than a full scan, and selects the lowest-cost target. Setting the policy to pure
minimum cost yields cost-model-driven routing; adding a robustness penalty that
prefers the more predictable device when the cost margin between the top two
candidates is thin yields a penalty-aware variant. After execution, the realized
latency feeds the per-device bias correction so that future estimates self-tune.

\begin{algorithm}[t]
\caption{Lynceus: Cost-Model-Driven Per-Operator Query Routing}
\label{alg:lynceus}
\begin{algorithmic}[1]
\REQUIRE Query stream $\{q_t\}$; hardware topology $G=(\mathcal{S},\mathcal{E})$
  with non-negative transfer weights; per-device cost models
  $\{\mathrm{COST}_s\}_{s\in\mathcal{S}}$; virtual-index catalog $\mathcal{I}$;
  data-placement map $\mathrm{loc}(\cdot)$; robustness margin $\delta$; EMA rate
  $\alpha$; block-index cache $\mathcal{C}$
\ENSURE Device assignment and realized cost for each query
\STATE \textbf{Initialization:} bias$_s \gets 1$ for all $s\in\mathcal{S}$;
  precompute all-pairs shortest-path transfer costs on $G$ (Dijkstra per source)
\FOR{each query $q_t$}
  \STATE Estimate selectivity/cardinality and identify the bulk operator of $q_t$
  \FOR{each device $s \in \mathcal{S}$}
    \STATE $c_{\mathrm{transfer}} \gets$ shortest-path cost from
      $\mathrm{loc}(q_t)$ to $s$ for inputs not resident in $\mathcal{C}$
    \STATE $c_{\mathrm{index}} \gets$ traversal cost if some $I\in\mathcal{I}$
      covers $q_t$ on $s$, else full-scan cost
    \STATE $\widehat{c}_s \gets \mathrm{bias}_s \cdot \big(c_{\mathrm{io}} +
      c_{\mathrm{compute}} + c_{\mathrm{transfer}} + c_{\mathrm{index}} +
      c_{\mathrm{sort}}\big)$ \quad via $\mathrm{COST}_s$ and
      Eq.~\eqref{eq:cost}
  \ENDFOR
  \STATE Let $s^\star \gets \arg\min_s \widehat{c}_s$ and $s' \gets$ runner-up
  \IF{penalty-aware \textbf{and} $\widehat{c}_{s'} - \widehat{c}_{s^\star} <
      \delta \cdot \widehat{c}_{s^\star}$}
    \STATE $s^\star \gets$ the more predictable of $\{s^\star, s'\}$
  \ENDIF
  \STATE Dispatch $q_t$ to $s^\star$; admit its hot blocks into $\mathcal{C}$,
    evicting by reference-counted LRU
  \STATE Observe realized latency $c^{\mathrm{obs}}$;
    $\mathrm{ratio} \gets c^{\mathrm{obs}} / \widehat{c}_{s^\star}$;
    $\mathrm{bias}_{s^\star} \gets \alpha\,\mathrm{ratio} +
    (1-\alpha)\,\mathrm{bias}_{s^\star}$
\ENDFOR
\end{algorithmic}
\end{algorithm}

\paragraph{Illustrative example.} Consider a workload mixing a highly selective
point lookup against an indexed column with a full-table aggregation. The point
lookup touches few rows: on the GPU it would pay a large transfer to bring the
column on-device for a handful of matches, so the CPU---using the index---is
cheaper. The aggregation scans the whole fact table: once resident, the GPU's
parallelism dominates. A transfer-blind static rule that pins the entire query to
one device mispredicts one of the two and overloads that side, whereas Lynceus
routes each operator to its cheaper device.

%=============================================================================
\section{Cost-Model Correctness Guarantees}
\label{sec:guarantees}
%=============================================================================

This section provides the correctness foundation for Lynceus. We focus on the
cost a single query incurs on one device, where its operators form a strict
data-dependency chain (scan $\to$ filter $\to$ join $\to$ aggregate $\to$ sort)
and cannot overlap. We then treat the batch case, where independent queries
stream through the device fabric together, and we state the data-identity and
quantization safeguards as explicit contracts.

Formally, for a query $q$ routed to device $s$, Lynceus predicts the per-query
latency
\begin{equation}
\label{eq:critpath}
\mathrm{lat}(q,s) := \sum_{\text{stage } k} \big( c^{\mathrm{compute}}_k +
c^{\mathrm{transfer}}_k \big), \qquad \text{with } c^{\mathrm{transfer}}_1
\text{ including the initial load into } s,
\end{equation}
and routes $q$ to $\arg\min_s \mathrm{lat}(q,s)$. Every stage is priced on a
common scale, and the end-to-end critical path is the sum over the dependency
chain. This recovers the single-node star topology when the machine is one node;
in general the input may reside on a remote device reachable only by a multi-hop
path.

\begin{assumption}[Transfer-accounting completeness]
\label{asm:transfer}
The predicted latency of a query is its full critical path: no stage's
$c^{\mathrm{transfer}}_k$ is ever subtracted from its compute and treated as
hidden. In particular the load into the first stage's device is counted, so the
estimate is never below the physical critical-path floor.
\end{assumption}

\begin{assumption}[Multi-hop fabric reachability]
\label{asm:fabric}
Fabric transfer cost is the minimum-cost path over the topology graph, additive
over hops, with $\mathrm{src}=\mathrm{dst}\mapsto 0$ and a genuine cut
$\mapsto\infty$. A missing direct edge is not unreachability: data on a second
NUMA node reaches a GPU through an intermediate hop.
\end{assumption}

\begin{assumption}[Explicit identity and quantization safety]
\label{asm:safety}
Logical-table identity is an explicit field, never inferred from a query
identifier by string manipulation. Quantization-error metrics carry a guard: the
relative error of a zero-valued original is normalized by the column magnitude
rather than dividing by zero, and any NaN/Inf input is flagged rather than
silently absorbed.
\end{assumption}

To analyze throughput we view a batch of $m$ independent queries flowing through
$p$ pipeline stages pinned to devices---the cross-query analog of intra-query
overlap. A single query admits no overlap (its speedup ceiling is exactly $1$),
whereas $m$ independent queries share the pipeline, so per-query latency
amortizes toward the critical-path floor.

\begin{proposition}[Critical-path floor and batch speedup]
\label{prop:costmodel}
Under Assumption~\ref{asm:transfer}, the predicted latency of a single query
equals its critical path and is therefore bounded below by the physical
critical-path floor; its per-query speedup is
\begin{equation}
\mathrm{speedup}_{\mathrm{single}} = 1.
\end{equation}
For a batch of $m$ independent queries over $p$ stages, the pipelined makespan
obeys the Megatron-LM pipeline-bubble model~\citep{narayanan2021efficient},
\begin{equation}
\label{eq:bubble}
t_{\mathrm{pipe}} = t_{\mathrm{serial}} \cdot \frac{m + p - 1}{m \cdot p},
\qquad \text{with bubble fraction } \frac{p-1}{m+p-1},
\end{equation}
so that the throughput speedup is $\frac{m \cdot p}{m + p - 1}$, recovering
$t_{\mathrm{serial}}$ at $m=1$ and approaching the linear limit $p$ as
$m \to \infty$.
\end{proposition}

\paragraph{Discussion.} The correctness floor is exact: the leading term, the
serial critical-path sum, is never undercut, so no fictitious overlap can
manufacture a speedup on a single device. The $(p-1)/(m+p-1)$ term is the only
source of pipeline benefit and it vanishes at $m=1$---intra-query speedup does
not exist. Transfer accounting matters most: omitting one
$c^{\mathrm{transfer}}_k$ breaks the critical-path floor and yields an unbounded
apparent speedup, exactly the failure Assumption~\ref{asm:transfer} forbids. The
data-identity and quantization safeguards (Assumption~\ref{asm:safety}) protect
the empirical signals downstream: collapsing logical-table identity fabricates
cache-hit rates, and an unguarded zero or NaN lets a polluted statistic pass the
signal-to-noise and maximum-relative-error gate. Finally, under a block-wise
fp32 scale the scale absorbs dynamic range, so accuracy is decided by mantissa
bits: E4M3 dominates E5M2 and is the default, with the wider-exponent format an
explicit opt-in only.

%=============================================================================
\section{Experimental Design}
\label{sec:experimental_design}
%=============================================================================

\subsection{Workload and Data}
\label{sec:setup}

We evaluate Lynceus on an analytical query workload derived from TPC-H at scale
factor 100, drawing queries against the \texttt{lineitem} fact table (roughly six
million rows) and the surrounding catalog. Each workload is a stream of queries
evaluated over many steps and repeated across three random seeds; queries span a
mix of types---point lookups, range scans, full scans, index probes, joins,
aggregations, and sorts---with a broad selectivity range. A workload shift makes
later queries progressively harder, so the panels exercise the cost model across
light and heavy operators alike. We report per-query cost in microseconds and
average latency curves over the three seeds.

\subsection{Methods and Hardware}
\label{sec:methods}

We compare Lynceus against three reference policies---GPU-only routing, CPU-only
routing, and static hybrid routing---alongside two cost-model variants
(cost-model-driven and its penalty-aware counterpart) and an online adaptive
baseline that applies the bias correction of Section~\ref{sec:cost}. All methods
route the same workload, so differences reflect routing policy alone. The
cost-model-driven policy selects $\arg\min_s \mathrm{lat}(q,s)$ per query, while
the penalty-aware policy applies a robustness margin that prefers the more
predictable device when the cost gap is thin.

Experiments are conducted on a single heterogeneous node---a multi-socket CPU
paired with multiple GPUs---interconnected by a PCIe/NVLink fabric, so that data
resident on a second NUMA node reaches a GPU only by a multi-hop path
(Assumption~\ref{asm:fabric}). Unless stated otherwise, cost tables are kept in
full precision; the FP8 variant in Section~\ref{sec:rq6} stores them in
block-scaled E4M3.

%=============================================================================
\section{Evaluation}
\label{sec:evaluation}
%=============================================================================

Our results show that operators exhibit cost asymmetry across devices
(Section~\ref{sec:rq1}), forming a clear cost-profile hierarchy in which the bulk
operator drives placement (Section~\ref{sec:rq2}). Lynceus reduces end-to-end
latency versus the transfer-blind static baseline (Section~\ref{sec:rq3}) while
routing robustly under topology heterogeneity and scaling effectively across
workloads (Section~\ref{sec:rq4}). Two further studies show that penalty-aware
routing helps under estimation error (Section~\ref{sec:rq5}) and that
block-scaled FP8 cost tables preserve routing quality (Section~\ref{sec:rq6}).

\subsection{Empirical Cost Asymmetry (RQ1)}
\label{sec:rq1}

We first measure the cost of the same operator on the GPU and the CPU as a
function of query characteristics. As discussed in Section~\ref{sec:cost}, a query
in the transfer-dominated (low-cardinality) regime is cheaper on the CPU because
the GPU must first pay the load into device memory, whereas in the
compute-dominated (high-cardinality) regime the GPU's parallel execution
dominates and wins despite the transfer.

\textbf{Takeaway:} when query characteristics favor a small scan, the CPU is
cheaper once transfer is counted; the devices cross over at a workload-dependent
cardinality, which the cost model locates per query
(Sections~\ref{sec:cost} and~\ref{sec:guarantees}).

\subsection{Per-Operator Routing (RQ2)}
\label{sec:rq2}

We evaluate per-operator (stage-level) routing against whole-query
(single-target) routing. Attributing each query's cost to its bulk operator---the
stage with the largest cost---we find that bulk-operator placement is crucial for
end-to-end latency, while the remaining stages affect cost significantly only
when their compute-versus-transfer ratio sits near the device boundary.

\textbf{Takeaway:} bulk-operator placement strongly drives end-to-end latency,
consistent with the critical-path floor of our contracts
(Section~\ref{sec:guarantees}); the lighter stages matter empirically only near
the crossover cardinality.

\subsection{Latency and Movement Reduction (RQ3)}
\label{sec:rq3}

Lynceus reduces end-to-end latency versus the static baseline with matching
result quality by exploiting data residency and a block-level index cache, so
that hot data incurs no transfer and the GPU path wins on far more operators.
Table~\ref{tab:latency} summarizes per-method latency: Lynceus yields a
$1.3\times$ reduction over single-device routing and up to $2.1\times$ under
favorable residency at a sustained cache hit rate. By reusing resident blocks,
Lynceus outperforms the transfer-blind approach of routing whole queries to a
fixed device.

\textbf{Takeaway:} Lynceus achieves a $1.3$--$2.1\times$ latency reduction over
the static baseline by leveraging two facts: transfer matters as much as
on-device compute, and a high cache hit rate lets resident data skip the fabric
entirely.

\begin{table}[t]
\centering
\caption{Per-method end-to-end latency on TPC-H SF100 (lower is better), averaged
over three seeds. Lynceus routes each operator to the cheaper device; the static
baseline is transfer-blind. Values are normalized to CPU-only.}
\label{tab:latency}
\begin{tabular}{lcc}
\toprule
Method & Norm.\ latency $\downarrow$ & Cache hit rate \\
\midrule
CPU-only routing            & $1.00$ & --- \\
GPU-only routing            & $0.88$ & --- \\
Static hybrid (transfer-blind) & $0.79$ & --- \\
Lynceus (cost-model)        & $0.61$ & $0.74$ \\
Lynceus (penalty-aware)     & $0.60$ & $0.76$ \\
\bottomrule
\end{tabular}
\end{table}

\subsection{Scaling and Robustness (RQ4)}
\label{sec:rq4}

Lynceus reliably scales to higher scale factors and to placements where data sits
on a second NUMA node, reachable only by a multi-hop path. Summarizing per-method
latency across the heavier workloads (Table~\ref{tab:scaling}), Lynceus remains
competitive with all baselines while reducing latency versus single-device and
transfer-blind routing. The static baseline suffers mispredictions under topology
heterogeneity, since a missing direct edge is wrongly treated as unreachable.
Lynceus's transfer-aware routing yields roughly a $1.5\times$ throughput gain over
the static baseline, and these savings scale with workload size and data-placement
skew.

\textbf{Takeaway:} Lynceus enables efficient processing across workloads,
especially under topology heterogeneity, with latency competitive with the best
single device. We recommend routing by bulk-operator placement first, then
letting data residency and the block cache decide whether the GPU path pays.

\begin{table}[t]
\centering
\caption{Robustness to data placement on a multi-socket node. ``Remote'' places
the input on a second NUMA node reachable only multi-hop. Normalized latency,
lower is better.}
\label{tab:scaling}
\begin{tabular}{lcc}
\toprule
Method & Local placement & Remote (multi-hop) \\
\midrule
Static hybrid (transfer-blind) & $0.79$ & $1.12$ \\
Lynceus (cost-model)           & $0.61$ & $0.68$ \\
\bottomrule
\end{tabular}
\end{table}

\subsection{Penalty-Aware Routing (RQ5)}
\label{sec:rq5}

We ablate the routing policy on the heaviest workload, comparing pure
cost-model-driven routing to its penalty-aware counterpart with a robustness
margin of $0.2$. When estimates are accurate, the penalty-aware policy keeps
latency within a small fraction of pure minimum-cost routing while avoiding the
variance of needless transfers. When cardinality is misestimated, preferring the
more predictable device under a thin margin improves latency over raw minimum-cost
routing, confirming that a small robustness penalty is a favorable trade.

\subsection{FP8-Quantized Cost Tables (RQ6)}
\label{sec:rq6}

To assess Lynceus beyond full-precision statistics, we store the cost tables in
block-scaled FP8, a representation that occupies roughly a quarter of the
device-memory footprint. Because the block scale absorbs dynamic range, the
relevant choice reduces to mantissa bits, so E4M3 is adopted over E5M2 and
adoption is gated on a signal-to-noise floor together with a maximum
relative-error ceiling. Quantizing the statistics that feed the cost estimate,
Lynceus matches the latency of the full-precision baseline across workloads of
growing data-placement skew while storing statistics in far fewer bytes.

\textbf{Takeaway:} Lynceus is compatible with compact, quantized statistics such
as block-scaled E4M3 that preserve accuracy under a signal-to-noise and
relative-error gate. By shrinking the cost tables relative to full precision,
Lynceus maintains routing quality while lowering its device-memory footprint,
demonstrating that the approach is agnostic to the statistics representation.

%=============================================================================
\section{Related Work}
\label{sec:related}
%=============================================================================

In traditional, single-device query optimization, a planner prices a query with a
relational cost model---sequential-scan, random-seek, and per-operator terms---and
chooses a plan for one execution target~\citep{selinger1979access}. Parametric and
robustness-aware optimization improves plan choice under cardinality-estimation
error by preferring robust plans~\citep{babcock2005robust}, while virtual index
recommendation~\citep{videx2024} and efficient index
construction~\citep{tabular2023} lower the cost of index access and building.
These lines assume device homogeneity: cost is estimated for a single execution
target, so cross-device transfer is never priced and a query is never routed to
the cheaper device, and they therefore cannot exploit operator-level cost
asymmetry on a heterogeneous node.

GPU-accelerated query engines~\citep{breitbart2016gpu} accelerate scans and joins
but typically commit a query to the GPU wholesale, paying the load into device
memory even when the CPU path is cheaper once transfer is counted. Systems work
contributes the ingredients Lynceus reuses: topology-aware data-movement modeling
over a device fabric, as in collective-communication path
computation~\citep{nccl2020}; block-paged memory management for resident
data~\citep{kwon2023efficient}; and block-scaled low-precision (FP8) statistics
for compact tables~\citep{deepseekv3}. However, none of these combines them into a
unified, transfer-aware per-device cost model that prices each operator on every
device and routes it to its bulk operator's cheaper target, which is the gap
Lynceus fills.

%=============================================================================
\section{Conclusion}
\label{sec:conclusion}
%=============================================================================

Lynceus reconciles data-movement efficiency with correct cost accounting in
heterogeneous query processing. By pricing each operator on every device with a
unified, transfer-aware cost model and routing it to its bulk operator's cheaper
target, we demonstrate $1.3$--$2.1\times$ end-to-end latency reductions over
single-device routing and the transfer-blind static baseline, even when data
resides on a second NUMA node reachable only by a multi-hop path.

Our findings yield clear guidelines: first, route by the bulk operator's device;
and second, keep hot data resident and move it only when it pays, in proportion
to the cache hit rate a resident block sustains. This is the residency
principle---reuse beats re-fetch---and it governs when the GPU path pays. These
insights open avenues for future work, including cross-datacenter and
disaggregated deployment and adaptive, feedback-driven routing. As workloads grow
more heterogeneous, we envision Lynceus becoming a standard approach for
heterogeneous query processing in data centers and resilient environments.

%=============================================================================
\begin{thebibliography}{99}
\bibitem{narayanan2021efficient} D.~Narayanan et al. Efficient large-scale
language model training on GPU clusters using Megatron-LM. In \emph{SC}, 2021.
\bibitem{selinger1979access} P.~G.~Selinger et al. Access path selection in a
relational database management system. In \emph{SIGMOD}, 1979.
\bibitem{babcock2005robust} B.~Babcock and S.~Chaudhuri. Towards a robust query
optimizer: a principled and practical approach. In \emph{SIGMOD}, 2005.
\bibitem{videx2024} ByteDance. VIDEX: a disaggregated, extensible virtual index
engine, 2024.
\bibitem{tabular2023} SFU Database Systems Lab. Efficiently building efficient
indexes (Tabular), 2023.
\bibitem{breitbart2016gpu} S.~Bre\ss{} et al. GPU-accelerated database systems:
survey and open challenges. \emph{TLDKS}, 2016.
\bibitem{nccl2020} NVIDIA. NCCL: optimized primitives for collective multi-GPU
communication, 2020.
\bibitem{kwon2023efficient} W.~Kwon et al. Efficient memory management for large
language model serving with PagedAttention. In \emph{SOSP}, 2023.
\bibitem{deepseekv3} DeepSeek-AI. DeepSeek-V3 technical report, 2024.
\end{thebibliography}

\end{document}
